\section{Structure of the optimal disposition}\label{sec:structure}

All proofs are in Appendix~\ref{app:proofs}.

\subsection{Pricing a unit without enumerating its plans}

For $\mu\ge 0$ let $\varphi_i(\mu)=\max_{s\in\cS_i}(g_{is}-\mu h_{is})$ be the best \emph{reduced margin} of a unit
of type $i$ when an operator hour is priced at $\mu$ on top of its cost.

\begin{proposition}[Pricing by access level]\label{prop:pricing}
Under Assumption~\ref{as:nested}, for every $i$ and $\mu\ge0$, with $\omega_i=w_i+\mu$,
\begin{equation}\label{eq:pricing}
\varphi_i(\mu)\;=\;\max_{\ell\in\cL}\Big\{-\omega_i\,a_{i\ell}+\sum_{j\in\cJ_i}\ \max_{d:\,\ell_d\le\ell}\
n_{ij}\big(\rho_{ijd}-\omega_i\,t_{jd}\big)\Big\}.
\end{equation}
The pricing problem of a type is solved in $O(|\cJ_i|\,|\cD|\,L)$ operations.
\end{proposition}

The set-up couples the components of a unit, which is why plans cannot be priced component by component in general.
Proposition~\ref{prop:pricing} says that the coupling has only $L+1$ states: once the access level is fixed,
components are independent. Assumption~\ref{as:nested} is needed: when a rough opening takes longer than a careful
one, the right-hand side of \eqref{eq:pricing} can overstate $\varphi_i$.

\begin{remark}[Two levels, without Assumption~\ref{as:nested}]\label{rem:twolevel}
In the instance of Section~\ref{sec:model}, write $\Psi^{d}_j=n_{ij}(\rho_{ijd}-\omega_it_{jd})$ for
$d\in\{\R,\Sd\}$. Then
\[
\varphi_i(\mu)=\max\Big\{0,\ \ -\omega_ia_{i1}+\sum_j\pos{\Psi^{\Sd}_j},\ \
-\omega_ia_{i2}+\sum_j\max(\Psi^{\R}_j,\Psi^{\Sd}_j,0)-\min_j\pos{\max(\Psi^{\Sd}_j,0)-\Psi^{\R}_j}\Big\},
\]
where the second term is kept only if some $\Psi^{\Sd}_j\ge0$. The three terms are the best plans with no opening,
with sorting only, and with at least one reuse; in the last, each component takes its best disposition, and if none
chooses reuse it is imposed on the component for which it costs least. No ordering of $a_{i1}$ and $a_{i2}$ is
used. This is the routine of the implementation.
\end{remark}

\subsection{The value of an hour}

\begin{proposition}[Dual and value function]\label{prop:dual}
\begin{equation}\label{eq:dual}
V(N,H)\;=\;\min_{\mu\ge0}\Big\{\mu H+\sum_{i\in\cI}N_i\,\varphi_i(\mu)\Big\}.
\end{equation}
$V$ is concave, piecewise linear and positively homogeneous of degree one in $(N,H)$, and nondecreasing in each
argument. If $\mu^\ast$ attains the minimum, then $V(N,H')\le V(N,H)+\mu^\ast(H'-H)$ for every $H'\ge0$. Any optimal
$\mu^\ast$ lies between the value of one more operator hour and the loss from one hour less, which coincide except
at the finitely many breakpoints of $V$; and $\mu^\ast=0$ is optimal when the labour budget is not exhausted.
\end{proposition}

The dual objective is convex and piecewise linear in $\mu$, with subgradient $H-\sum_iN_ih_{is_i(\mu)}$ at any
selection $s_i(\mu)$ of best plans. The implementation locates $\mu^\ast$ by bisection on the sign of this
subgradient, each evaluation costing one pricing per type by Proposition~\ref{prop:pricing}, and recovers a primal
solution by Proposition~\ref{prop:split}.

\begin{proposition}[Few split types, efficient plans]\label{prop:split}
Problem \eqref{eq:P} has an optimal solution in which every type with $N_i>0$ uses one plan, except at most one
type, which uses two. In any optimal solution, every plan in use attains $\varphi_i(\mu^\ast)$ for an optimal
multiplier $\mu^\ast$; in particular it lies on the upper concave envelope of the points
$\{(h_{is},g_{is})\}_{s\in\cS_i}$, and the two plans of a split type lie on a common edge of that envelope, of
slope $\mu^\ast$. With $M$ coupling constraints in place of the single labour budget, a feasible programme has an
optimal solution in which at most $M$ types use more than one plan.
\end{proposition}

For a single budget this is the classical structure of the linear multiple-choice knapsack
\citep{sinha1979,zemel1980}; we state it for $M$ constraints because circularity targets add constraints.

\begin{remark}[Integrality]\label{rem:integrality}
With a single budget, rounding the split type down to whole units and sending the remainder to the cheaper of its
two plans keeps the solution feasible and loses less than the margin difference of one unit. The argument holds for
$\le$ constraints only; a lower bound on an outcome can be violated by rounding down.
\end{remark}

\paragraph{Managerial reading.} At the optimum a single number, the value of an hour $\mu^\ast$, coordinates all
decisions: every type is processed with the plan that is best when labour is charged $w_i+\mu^\ast$, and at most
one type is processed in two ways, because the hours run out part-way through it. The same number tells the manager
what an additional operator hour is worth, which is the first question asked once a line exists.

\subsection{A threshold rule for the bench}

Consider a unit that is opened carefully anyway. The choice is then separable across components.

\begin{proposition}[Yield threshold]\label{prop:threshold}
In the instance of Section~\ref{sec:model}, let $\omega=w_i+\mu$ and assume $\bar v_j>m_jp_j$. In a unit opened to
level 2, reuse is at least as good as both sorting and leaving component $j$ if and only if
\begin{equation}\label{eq:threshold}
q_{ij}\;\ge\;q^{\dagger}_{ij}(\mu)\;=\;\frac{\kappa_j+\omega\,t_{j\R}-\min\{\omega\,t_{j\Sd},\ m_j(p_j+\tau)\}}
{\bar v_j-m_jp_j}.
\end{equation}
When $t_{j\R}\ge t_{j\Sd}$ the threshold is nondecreasing in $\mu$.
\end{proposition}

The rule is what an operator can apply: \emph{extract and test component $j$ when its observed good rate is above
$q^\dagger_{ij}$}. It also says how scarcity of labour changes behaviour: the scarcer the hours, the higher the
yield a component must show. Both hypotheses are needed: the comparison reverses when a good copy is worth less
reused than as material, and the threshold can decrease in $\mu$ when testing is faster than sorting. The threshold
is conditional on the unit being opened carefully; whether it is opened at all is decided by
Proposition~\ref{prop:pricing}, where the set-up is weighed against the sum of the components' reduced values.

\subsection{Circularity objectives}

Three outcomes are linear in $x$ and are reported for every solution: the mass diverted from the bulk stream,
$\sum_{i,s}b_{is}x_{is}$ with $b_{is}$ the mass taken out of a unit under plan $s$; the mass of critical raw
materials sent to a recovery route; and the number of parts returned to service, $\sum_{i,s}e_{is}x_{is}$ with
$e_{is}=\sum_{j:s_j=\R}n_{ij}q_{ij}$.

\begin{corollary}[Frontier]\label{cor:frontier}
Adding a weighted linear outcome to the objective keeps at most one split type. Imposing $M$ lower bounds on linear
outcomes gives $M+1$ coupling constraints; when the programme is feasible it has an optimal solution with at most
$M+1$ split types.
\end{corollary}

The frontier between margin and parts returned to service answers a question that a single-objective optimum
cannot: how much margin does each additional part returned to service cost? Returning parts to service lowers the
cost of repair and of refurbished systems; whether this improves access to imaging for hospitals with small budgets
is an empirical question outside the model, and we make no claim about it.
